\documentclass[11pt]{amsart}
\usepackage[margin=1.1in]{geometry}
\usepackage{amsmath,amssymb,amsthm,mathtools}
\usepackage{enumitem}
\usepackage{microtype}
\usepackage{hyperref}
\usepackage[nameinlink,capitalize]{cleveref}
\hypersetup{colorlinks=true,linkcolor=blue,citecolor=blue,urlcolor=blue}
\pdfstringdefDisableCommands{%
  \def\({}%
  \def\){}%
  \def\tilde#1{#1}%
  \def\mathfrak#1{#1}%
  \def\mathrm#1{#1}%
  \def\mathbb#1{#1}%
  \def\tau{tau}%
  \def\lambda{lambda}%
  \def\infty{infty}%
}

\theoremstyle{plain}
\newtheorem{theorem}{Theorem}[section]
\newtheorem{lemma}[theorem]{Lemma}
\newtheorem{proposition}[theorem]{Proposition}
\newtheorem{corollary}[theorem]{Corollary}
\theoremstyle{definition}
\newtheorem{definition}[theorem]{Definition}
\newtheorem{assumption}[theorem]{Assumption}
\theoremstyle{remark}
\newtheorem{remark}[theorem]{Remark}

\title[Scale Collapse for Type II Rescalings]{Scale-Collapse Cost Exclusion for Type II Navier--Stokes Rescalings}
\author{Guillem Cordero}
\date{}

\begin{document}

\begin{abstract}
We work with suitable weak solutions of the three-dimensional Navier--Stokes equations and give an explicit step-by-step version of the scale-collapse exclusions.  Under a renormalized gauge, the logarithmic scale satisfies
\(\partial_\tau\log\lambda=a\).  If a renormalized branch has a nontrivial localized core and the scale collapses, then
\(\int a_-(\tau)M_R(\tau)\,d\tau=\infty\); thus finite scale-collapsing cost or finite absolute cost is incompatible with collapse.  In addition, we record the autonomous-window limit passage under an explicitly stated compactness input; that input is not derived from the cost hypotheses in this paper.
\end{abstract}

\maketitle

\paragraph{Notation and functional conventions.}
For a bounded spatial region \(K\subset\mathbb R^3\), we use
\[
  (f,g)_{L^2(K)}:=\int_K f\cdot g\,dx,\qquad
  \langle F,\phi\rangle_{H^{-1}(K),H^1_0(K)}
  :=\int_0^T\langle F(\cdot,s),\phi(\cdot,s)\rangle_{H^{-1}(K),H^1_0(K)}\,ds
\]
whenever \(F\in L^1(0,T;H^{-1}(K))\) and \(\phi\in L^\infty(0,T;H^1_0(K))\). Weak limits and weak convergences are taken in the distributional sense on \(C_c^\infty\)-test functions.
\[
  \|f\|_{L^p_tL^q_x}:=\left(\int_0^T\|f(\cdot,s)\|_{L^q(\mathbb R^3)}^p\,ds\right)^{1/p},
\]
for any fixed \(T>0\), with the usual modification \(p=\infty\).

\section{Renormalized scale-collapse setting}

\paragraph{Position of the result.}
We split the paper into two separate levels:
\begin{enumerate}[label=(\arabic*)]
  \item Part A (Sections~\ref{sec:cost}--\ref{sec:finite-cost}): unconditional cost
    exclusion arguments that do not use compactness.
  \item Part B (Section~\ref{sec:autonomous}): a conditional compactness statement under explicitly
    stated compactness hypotheses.
\end{enumerate}

Throughout this paper, we assume the renormalized representation established earlier: on the
time interval under consideration, the fields \((V,P)\) satisfy \eqref{eq:renorm-ns} and its
distributional weak formulation.  In particular, the solenoidal formulation is obtained by
restricting the test fields below to divergence-free fields.

Let \((u,p)\) be a divergence-free suitable weak solution of
\[
  \partial_tu+(u\cdot\nabla)u+\nabla p=\nu\Delta u,
  \qquad \nabla\cdot u=0,
\]
in \(\mathbb R^3\), with viscosity \(\nu>0\).  Let
\(\lambda(\tau)>0\), \(X(\tau)\in\mathbb R^3\), and a measurable gauge function
\(c(t)\) be given, and define the renormalized fields by
\[
  u(x,t)=\lambda(\tau)^{-1}V\left(\frac{x-X(\tau)}{\lambda(\tau)},\tau\right),
  \qquad
  p(x,t)=\lambda(\tau)^{-2}P\left(\frac{x-X(\tau)}{\lambda(\tau)},\tau\right)+c(t).
\]
Pressure is defined only up to an arbitrary function of time; \(c(t)\) has no effect on
\(\nabla p\), so it is irrelevant for all subsequent estimates.
Assume that \((V,P)\) satisfies on \((0,\infty)\times\mathbb R^3\)
\begin{equation}\label{eq:renorm-ns}
  \partial_\tau V+(V\cdot\nabla)V+\nabla P
  =\nu\Delta V+a(\tau)(V+y\cdot\nabla V)+b(\tau)\cdot\nabla V,
  \qquad \nabla\cdot V=0.
\end{equation}
For every \(\psi\in C_c^\infty((0,\infty)\times\mathbb R^3;\mathbb R^3)\) with \(\nabla\cdot\psi=0\),
the pair \((V,P)\) satisfies the weak formulation
\[
\begin{aligned}
0={}&\int_0^\infty\!\!\int_{\mathbb R^3}
\bigl(
-V\cdot\partial_\tau\psi
-(V\otimes V):\nabla\psi
+\nu\nabla V:\nabla\psi
\bigr)\,dy\,d\tau \\
&+\int_0^\infty\!\!\int_{\mathbb R^3}
\bigl(
-a(\tau)(V+y\cdot\nabla V)\cdot\psi
-(b(\tau)\cdot\nabla V)\cdot\psi
+P\,\nabla\cdot\psi
\bigr)\,dy\,d\tau .
\end{aligned}
\]
For every \(0<T<\infty\), restricting the \(\tau\)-integration to \((0,T)\) gives the same identity on \((0,T)\times\mathbb R^3\).
Equivalently, for every
\(\psi\in C_c^\infty((0,\infty)\times\mathbb R^3;\mathbb R^3)\), not necessarily
divergence free,
\[
\begin{aligned}
0={}&\int_0^\infty\!\!\int_{\mathbb R^3}
\bigl(
-V\cdot\partial_\tau\psi
-(V\otimes V):\nabla\psi
\nu\nabla V:\nabla\psi
-P\,\nabla\cdot\psi
\bigr)\,dy\,d\tau \\
&+\int_0^\infty\!\!\int_{\mathbb R^3}
\bigl(
-a(\tau)(V+y\cdot\nabla V)\cdot\psi
-(b(\tau)\cdot\nabla V)\cdot\psi
\bigr)\,dy\,d\tau .
\end{aligned}
\]

Define positive and negative parts of the scalar field \(a\) by
\[
  a_+(\tau):=\max\{a(\tau),0\},
  \qquad
  a_-(\tau):=\max\{-a(\tau),0\}.
\]
Then for a.e. \(\tau\), \(a=a_+-a_-\), \(a_+,a_-\ge0\), and \(a_+a_-=0\).  Assume
\[
  a\in L^1_{\mathrm{loc}}([\tau_0,\infty)),
  \qquad
  b\in L^1_{\mathrm{loc}}([\tau_0,\infty);\mathbb R^3),
\]
and the scale convention
\begin{equation}\label{eq:log-scale-convention}
  \frac{d}{d\tau}\log\lambda(\tau)=a(\tau)
\end{equation}
for a.e. \(\tau\ge\tau_0\), with \(\log\lambda\) absolutely continuous on compact intervals.

All integrals in this paper are interpreted in \([0,\infty]\). When we write
\(\int_{\tau_0}^{\infty}\), partial integrals are taken over \([\tau_0,T]\) and
monotone convergence is used as \(T\to\infty\).

Fix \(R>0\).  Choose
\(
  \Phi\in C_c^\infty(\mathbb R^3;[0,1])
\)
satisfying
\begin{equation}\label{eq:Phi-properties}
  \Phi(y)=1\ \text{if }|y|\le R,
  \quad
  \operatorname{supp}_y\Phi\subset B_{2R}.
\end{equation}
Define the localized kinetic mass
\[
  M_R(\tau):=\int_{\mathbb R^3}|V(y,\tau)|^2\Phi(y)
  \,dy.
\]
Since \(\Phi\) is fixed in \(\tau\), \(M_R\) is measurable in \(\tau\), and only this
measurability is used below.

\begin{definition}[Nonvanishing localized core]
A renormalized branch has nonvanishing localized core if there exist
constants \(m_0>0\) and \(\tau_1\ge\tau_0\) such that
\[
  M_R(\tau)\ge m_0\quad\text{for a.e. }\tau\in[\tau_1,\infty).
\]
In applications, this lower bound is supplied by the selected-core construction; in this
paper it is stated as an explicit hypothesis.
\end{definition}

\begin{definition}[Scale-collapse and absolute scale-drift densities]
Define
\[
  \mathcal C_{\mathrm{sc}}(\tau):=a_-(\tau)M_R(\tau),
  \qquad
  \mathcal C_{\mathrm{abs}}(\tau):=\nu\int_{\mathbb R^3}|\nabla V(y,\tau)|^2\Phi(y)
  \,dy+|a(\tau)|M_R(\tau).
\]
\(\mathcal C_{\mathrm{sc}}\) is the negative drift weighted by localized core mass, while
\(\mathcal C_{\mathrm{abs}}\) controls both collapse drift and local dissipation.
Their total costs are extended real-valued by
\[
  \int_{\tau_0}^{\infty}\mathcal C_{\mathrm{sc}}(\tau)
  \,d\tau,
  \qquad
  \int_{\tau_0}^{\infty}\mathcal C_{\mathrm{abs}}(\tau)
  \,d\tau.
\]
\end{definition}

\begin{definition}[Finite-cost branch]\label{def:finite-cost-branch}
A branch has finite scale-collapsing cost if
\[
  \int_{\tau_0}^{\infty}\mathcal C_{\mathrm{sc}}(\tau)
  \,d\tau<\infty.
\]
It has finite absolute cost if
\[
  \int_{\tau_0}^{\infty}\mathcal C_{\mathrm{abs}}(\tau)
  \,d\tau<\infty.
\]
\end{definition}

\section{Elementary alternative}
\label{sec:cost}

\begin{lemma}[Finite-or-infinite alternative]\label{lem:finite-infinite}
Let \(f\ge0\) be measurable on \([\tau_0,\infty)\). Then exactly one of the following holds:
\[
  \int_{\tau_0}^{\infty}f(\tau)
  \,d\tau<\infty
  \quad\text{or}\quad
  \int_{\tau_0}^{\infty}f(\tau)
  \,d\tau=\infty.
\]
\end{lemma}

\begin{proof}
\begin{enumerate}[label=(\arabic*)]
  \item For \(T>\tau_0\), define partial costs
  \[
    F(T):=\int_{\tau_0}^{T}f(\tau)
    \,d\tau\in[0,\infty].
  \]

  \item If \(T_2>T_1\), then by nonnegativity of \(f\),
  \[
    F(T_2)-F(T_1)=\int_{T_1}^{T_2}f(\tau)
    \,d\tau\ge0,
  \]
  so \(F\) is monotone nondecreasing.

  \item A monotone function on \([\tau_0,\infty)\) has an extended limit
  \(L\in[0,\infty]\),
  \[
    L:=\lim_{T\to\infty}F(T).
  \]
  By definition of improper integral, this limit equals
  \(\int_{\tau_0}^{\infty}f(\tau)\,d\tau\).

  \item Therefore \(L\) is either finite or equal to \(+\infty\), and these
  cases are mutually exclusive.
\end{enumerate}
Hence exactly one item in the lemma holds.
\end{proof}

\begin{lemma}[Logarithmic scale identity]\label{lem:log-scale}
Under \eqref{eq:log-scale-convention}, for all \(\tau_1,\tau_2\) with
\(\tau_0\le\tau_1<\tau_2<\infty\),
\[
  \int_{\tau_1}^{\tau_2}a(\tau)
  \,d\tau
  =\log\lambda(\tau_2)-\log\lambda(\tau_1).
\]
\end{lemma}

\begin{proof}
\begin{enumerate}[label=(\arabic*)]
  \item Since \(a\in L^1_{\mathrm{loc}}\), the map
  \(\log\lambda\) is absolutely continuous on \([\tau_1,\tau_2]\), hence has a
  representative with a.e. derivative equal to \(a\).

  \item By the fundamental theorem of calculus for absolutely continuous functions,
  for every \(\tau\in[\tau_1,\tau_2]\),
  \[
    \log\lambda(\tau)=\log\lambda(\tau_1)+\int_{\tau_1}^{\tau}a(s)
    \,ds.
  \]

  \item Setting \(\tau=\tau_2\) yields the identity.
\end{enumerate}
\end{proof}

\section{Scale collapse versus localized core lower bounds}

\begin{theorem}[Collapse with nonvanishing localized core forces infinite
scale-collapsing cost]\label{thm:collapse-core-floor}
Assume:
\begin{enumerate}[label=(\arabic*)]
  \item \(\lambda(\tau)\to0\) as \(\tau\to\infty\);
  \item the branch has nonvanishing localized core.
\end{enumerate}
Then
\[
  \int_{\tau_0}^{\infty}\mathcal C_{\mathrm{sc}}(\tau)
  \,d\tau
  =\int_{\tau_0}^{\infty}a_-(\tau)M_R(\tau)
  \,d\tau=\infty.
\]
\end{theorem}

\begin{proof}
\begin{enumerate}[label=(\arabic*)]
  \item Choose \(m_0>0\) and \(\tau_1\ge\tau_0\) from the nonvanishing
  core assumption so that
  \[
    M_R(\tau)\ge m_0\quad\text{for a.e. }\tau\in[\tau_1,\infty).
  \]

  \item By Lemma~\ref{lem:log-scale}, for any \(\tau_2>\tau_1\),
  \[
    \int_{\tau_1}^{\tau_2}a(\tau)
    \,d\tau
    =\log\lambda(\tau_2)-\log\lambda(\tau_1).
  \]
  Taking \(\tau_2\to\infty\) and using \(\lambda(\tau_2)\to0\),
  \[
    \lim_{\tau_2\to\infty}\int_{\tau_1}^{\tau_2}a(\tau)
    \,d\tau=-\infty.
  \]

  \item Decompose \(a=a_+-a_-\).  Then for each \(\tau_2>\tau_1\),
  \[
    \int_{\tau_1}^{\tau_2}a(\tau)
    \,d\tau
    =\int_{\tau_1}^{\tau_2}a_+(\tau)
    \,d\tau-\int_{\tau_1}^{\tau_2}a_-(\tau)
    \,d\tau
    \ge-\int_{\tau_1}^{\tau_2}a_-(\tau)
    \,d\tau.
    \]
    Hence
    \[
    \int_{\tau_1}^{\tau_2}a_-(\tau)
    \,d\tau\ge
    -\int_{\tau_1}^{\tau_2}a(\tau)
    \,d\tau.
    \]

  \item Taking limits and using Step~2 yields
  \[
    \int_{\tau_1}^{\infty}a_-(\tau)
    \,d\tau=\infty.
  \]
  Indeed, by Step~2,
  \[
    -\int_{\tau_1}^{\tau_2}a(\tau)\,d\tau\to+\infty.
  \]
  Since the partial integrals of \(a_-\) dominate this quantity, they are unbounded.
  Because \(a_-\ge0\), monotone convergence gives
  \(\int_{\tau_1}^{\infty}a_-(\tau)\,d\tau=\infty\).

  \item Split the cost into early and late times:
  \[
    \int_{\tau_0}^{\infty}a_-M_R
    \,d\tau
    =\int_{\tau_0}^{\tau_1}a_-M_R
    \,d\tau+
    \int_{\tau_1}^{\infty}a_-M_R
    \,d\tau.
  \]
  The first term is finite or infinite independent from the argument; the second term
  satisfies
  \[
    \int_{\tau_1}^{\infty}a_-(\tau)M_R(\tau)
    \,d\tau
    \ge m_0\int_{\tau_1}^{\infty}a_-(\tau)
    \,d\tau
    =\infty
  \]
  by the bound in (1) and Step~4.  Hence the total integral is infinite.
\end{enumerate}
\end{proof}

\section{Finite-cost exclusions}
\label{sec:finite-cost}

\begin{theorem}[Scale-collapse cost exclusion]\label{thm:cost-exclusion}
No branch with finite scale-collapsing cost (Definition~\ref{def:finite-cost-branch}) can satisfy simultaneously
\(\lambda(\tau)\to0\) and nonvanishing localized core.
\end{theorem}

\begin{proof}
Assume, for contradiction, that such a branch exists.  The two hypotheses are exactly
those of Theorem~\ref{thm:collapse-core-floor}; therefore
\[
  \int_{\tau_0}^{\infty}\mathcal C_{\mathrm{sc}}(\tau)
  \,d\tau=\infty.
\]
This contradicts finite scale-collapsing cost.  Therefore no such branch exists.
\end{proof}

\begin{lemma}[Absolute-cost domination]\label{lem:absolute-cost}
For a.e. \(\tau\in[\tau_0,\infty)\),
\[
  \mathcal C_{\mathrm{abs}}(\tau)\ge\mathcal C_{\mathrm{sc}}(\tau).
\]
Hence, pointwise and after integrating in \(\tau\),
\[
  \int_{\tau_0}^{\infty}\mathcal C_{\mathrm{sc}}(\tau)\,d\tau
  \le
  \int_{\tau_0}^{\infty}\mathcal C_{\mathrm{abs}}(\tau)\,d\tau,
\]
as extended real numbers.
In particular, finite absolute cost implies finite scale-collapsing cost, and
equivalently
\[
  \int_{\tau_0}^{\infty}\mathcal C_{\mathrm{sc}}(\tau)\,d\tau=\infty
  \quad\Longrightarrow\quad
  \int_{\tau_0}^{\infty}\mathcal C_{\mathrm{abs}}(\tau)\,d\tau=\infty.
\]
\end{lemma}

\begin{proof}
At each time \(\tau\),
\[
  \mathcal C_{\mathrm{abs}}(\tau)
  =\nu\int|\nabla V|^2\Phi
  +(a_++a_-)M_R
  \ge(a_++a_-)M_R\ge a_-M_R
  =\mathcal C_{\mathrm{sc}}(\tau),
\]
since \(a_+,a_-\ge0\) and \(\nu\int|\nabla V|^2\Phi\ge0\).  Integrating over
\([\tau_0,\infty)\) yields the claimed implication.
\end{proof}

\begin{corollary}[Absolute scale-cost exclusion]\label{cor:absolute-cost-exclusion}
No branch with finite absolute cost in the sense of Definition~\ref{def:finite-cost-branch}
can satisfy simultaneously \(\lambda(\tau)\to0\) and nonvanishing localized core.
\end{corollary}

\begin{proof}
Assume instead such a branch exists.  By Theorem~\ref{thm:collapse-core-floor},
\(\int\mathcal C_{\mathrm{sc}}=\infty\).  By Lemma~\ref{lem:absolute-cost}, the
absolute cost is pointwise above \(\mathcal C_{\mathrm{sc}}\), so its integral must also be
infinite, contradicting finite absolute cost.
\end{proof}

\begin{remark}
The exclusions in this section are purely integral and use neither compactness nor any
autonomous-limit argument.
\end{remark}

\section{Autonomous reduced-limit theorem from the selected-window estimates}
\label{sec:autonomous}

The exclusions above are integral statements and do not rely on compactness.  We now give the
standard local compactness argument on a fixed finite autonomous window.  The input used here is
not the cost bound alone; it is the selected-window suitable package supplied by the preceding
renormalized representation, pressure reconstruction, and local Caccioppoli estimates.

\paragraph{Dependencies used in this section.}
In standard PDE terms, the preceding papers supply the following ingredients on selected compact
renormalized windows:
\begin{enumerate}[label=(\arabic*)]
  \item the repaired-gauge renormalized equation and weak formulation for \((V,P)\);
  \item pressure reconstruction modulo spatial means, with local \(L^{3/2}\) bounds on compact
    cylinders;
  \item local windowed bounds
  \(V\in L^\infty_tL^2_x\cap L^2_tH^1_x\) on compact cylinders;
  \item a retained compact core lower bound preventing the selected profile from vanishing.
\end{enumerate}
Together with the thick-window convergence of the modulation coefficients, these estimates imply
the compactness and limit passage below.

\begin{definition}[Thick autonomous windows]\label{def:thick-autonomous-windows}
A branch has \emph{thick autonomous windows} if there exist constants
\(T>0\), \(c>0\), \(M_1\in(0,\infty)\), a sequence \((\tau_n)\) with
\(\tau_n\to\infty\), and limits \(a_\infty\in\mathbb R\),
\(b_\infty\in\mathbb R^3\), such that for every \(n\):
\begin{enumerate}[label=(\roman*)]
  \item
  \[
    \frac1T\int_{\tau_n}^{\tau_n+T}a_-(\tau)M_R(\tau)
    \,d\tau\ge c;
  \]
  \item
  \[
    0\le M_R(\tau)\le M_1\quad\text{for a.e. }\tau\in[\tau_n,\tau_n+T];
  \]
  \item
  \[
    a(\tau_n+\cdot)\to a_\infty\ \text{in }L^1(0,T),
    \qquad
    b(\tau_n+\cdot)\to b_\infty\ \text{in }L^1(0,T;\mathbb R^3).
  \]
\end{enumerate}
\end{definition}
\begin{remark}
This notion is used only to isolate a limiting drift coefficient \(a_\infty<0\).  The
compactness to extract a profile limit is obtained from the selected-window suitable estimates,
not from the scalar cost inequalities of Sections~\ref{sec:cost}--\ref{sec:finite-cost}.
\end{remark}

\begin{lemma}[Thick drift implies negative autonomous limit parameter]\label{lem:negative-parameter}
On thick autonomous windows, the limiting autonomous coefficient is negative:
\[
  a_\infty\le-\frac{c}{M_1}<0.
\]
\end{lemma}

\begin{proof}
\begin{enumerate}[label=(\arabic*)]
  \item From (ii), for a.e. \(\tau\in[\tau_n,\tau_n+T]\),
  \(0\le M_R(\tau)\le M_1\).  Therefore
  \[
    \frac{a_-(\tau)M_R(\tau)}{M_1}
    \le a_-(\tau).
  \]
  Therefore
  \[
    \frac1T\int_{\tau_n}^{\tau_n+T}a_-(\tau)
    \,d\tau
    \ge \frac1{M_1T}\int_{\tau_n}^{\tau_n+T}a_-(\tau)M_R(\tau)
    \,d\tau
    \ge \frac{c}{M_1}.
  \]

  \item Set \(a_n(s):=a(\tau_n+s)\).  By (iii), \(a_n\to a_\infty\) in
  \(L^1(0,T)\). The map \(r\mapsto r^-\) is 1-Lipschitz, so
  \(a_n^-\to a_\infty^-\) in \(L^1(0,T)\):
  \[
    \|a_n^--a_\infty^-\|_{L^1(0,T)}\le
    \|a_n-a_\infty\|_{L^1(0,T)}\to0.
  \]
  Thus
  \[
  \lim_{n\to\infty}\frac1T\int_{\tau_n}^{\tau_n+T}a_-(\tau)
    \,d\tau
    =\frac1T\int_0^T(a_\infty)^-
    \,ds=(a_\infty)^-.
  \]

  \item Combine Step~1 and Step~2:
  \((a_\infty)^-\ge c/M_1\), i.e.\ \(a_\infty\le-c/M_1<0\).
\end{enumerate}
\end{proof}

\begin{assumption}[Selected-window suitable estimates]\label{ass:selected-window-estimates}
Let \((\tau_n)\) be a thick autonomous-window sequence and set
\[
  V_n(y,s):=V(y,\tau_n+s),\qquad
  P_n(y,s):=P(y,\tau_n+s),
\]
\[
  \alpha_n(s):=a(\tau_n+s),\qquad
  \beta_n(s):=b(\tau_n+s),
  s\in(0,T).
\]
For every bounded Lipschitz domain \(K\Subset\mathbb R^3\), assume there is a constant
\(C_K<\infty\), independent of \(n\), such that
\[
  \|V_n\|_{L^\infty(0,T;L^2(K))}
  +\|V_n\|_{L^2(0,T;H^1(K))}
  +\|P_n-(P_n)_K\|_{L^{3/2}((0,T)\times K)}
  \le C_K,
\]
where
\[
  (P_n)_K(s):=\frac1{|K|}\int_KP_n(y,s)\,dy.
\]
Assume also that each \(V_n\) is divergence free and that there exist
\(\chi\in C_c^\infty(B_R)\), \(\chi\ge0\), \(\eta>0\), and \(N\in\mathbb N\) such that
for all \(n\ge N\),
\[
  \int_{\mathbb R^3}|V_n(y,s)|^2\chi(y)\,dy\ge\eta
  \quad\text{for a.e. }s\in(0,T).
\]
\end{assumption}

\begin{proposition}[Derivation of the selected-window estimates from the preceding papers]
\label{prop:selected-window-derived}
Let a repaired-gauge branch satisfy the renormalized representation and pressure reconstruction
of the preceding representation theorem on the windows \([\tau_n,\tau_n+T]\).  Assume that on
each compact cylinder \((0,T)\times K\) the local compact-cylinder Caffarelli--Kohn--Nirenberg
quantities are finite uniformly in \(n\), and that the selected core has the localized lower
bound stated in Assumption~\ref{ass:selected-window-estimates}.  Then
Assumption~\ref{ass:selected-window-estimates} holds.
\end{proposition}

\begin{proof}
The renormalized representation theorem supplies the shifted equation, divergence-free condition,
and pressure reconstruction modulo functions of time.  The pressure reconstruction gives, after
subtracting spatial means,
\[
  \sup_n\|P_n-(P_n)_K\|_{L^{3/2}((0,T)\times K)}<\infty
\]
on every compact \(K\).  The local Caccioppoli estimate applied on a slightly larger compact
cylinder gives the inner bounds
\[
  \sup_n\|V_n\|_{L^\infty(0,T;L^2(K))}
  +\sup_n\|V_n\|_{L^2(0,T;H^1(K))}<\infty.
\]
These are precisely the local velocity and pressure estimates in
Assumption~\ref{ass:selected-window-estimates}.  The retained selected-core lower bound is the
last condition in that assumption.  Hence all items in
Assumption~\ref{ass:selected-window-estimates} follow from the preceding representation,
pressure, and local Caccioppoli results.
\end{proof}

\begin{lemma}[Local time regularity on autonomous windows]\label{lem:window-time-regularity}
Assume thick autonomous windows and Assumption~\ref{ass:selected-window-estimates}.  Let
\(K_0\Subset K_1\Subset\mathbb R^3\) be bounded Lipschitz domains, and choose an integer
\(m\) so large that \(H^m_0(K_1)\hookrightarrow W^{1,\infty}_0(K_1)\).  Then
\[
  \{\partial_sV_n\}_{n\ge1}
  \quad\text{is bounded in }L^1(0,T;H^{-m}(K_1)).
\]
\end{lemma}

\begin{proof}
Fix \(\zeta\in C_c^\infty(K_1)\) with \(\zeta\equiv1\) on a neighborhood of \(K_0\).
Let \(\phi\in H^m_0(K_1;\mathbb R^3)\).  Since \(H^m_0(K_1)\hookrightarrow W^{1,\infty}_0(K_1)\),
there is \(C_m\) such that
\[
  \|\phi\|_{L^\infty(K_1)}+\|\nabla\phi\|_{L^\infty(K_1)}
  \le C_m\|\phi\|_{H^m(K_1)}.
\]
The shifted equation is
\[
  \partial_sV_n
  =-(V_n\cdot\nabla)V_n-\nabla P_n+\nu\Delta V_n
  +\alpha_n(V_n+y\cdot\nabla V_n)+\beta_n\cdot\nabla V_n
\]
in distributions.  We estimate each term against \(\phi\).

For the convection term, integration by parts gives
\[
  \left|\langle (V_n\cdot\nabla)V_n,\phi\rangle\right|
  =
  \left|\int_{K_1}V_{n,i}V_{n,j}\partial_j\phi_i\,dy\right|
  \le C_m\|V_n(s)\|_{L^2(K_1)}^2\|\phi\|_{H^m(K_1)}.
\]
For the viscous term,
\[
  \left|\langle \nu\Delta V_n,\phi\rangle\right|
  =
  \nu\left|\int_{K_1}\nabla V_n:\nabla\phi\,dy\right|
  \le \nu C_m\|\nabla V_n(s)\|_{L^2(K_1)}|K_1|^{1/2}\|\phi\|_{H^m(K_1)}.
\]
For the pressure term, subtracting the spatial mean does not change the gradient:
\[
  \left|\langle\nabla P_n,\phi\rangle\right|
  =
  \left|\int_{K_1}(P_n-(P_n)_{K_1})\,\nabla\cdot\phi\,dy\right|
  \le C_m |K_1|^{1/3}
  \|P_n-(P_n)_{K_1}\|_{L^{3/2}(K_1)}
  \|\phi\|_{H^m(K_1)}.
\]
For the \(V_n\)-part of the scaling drift,
\[
  |\alpha_n(s)|\,\left|\int_{K_1}V_n\cdot\phi\,dy\right|
  \le C_m|\alpha_n(s)|\,|K_1|^{1/2}
  \|V_n(s)\|_{L^2(K_1)}\|\phi\|_{H^m(K_1)}.
\]
For the \(y\cdot\nabla V_n\)-part, integrate by parts:
\[
  \int_{K_1}(y\cdot\nabla V_n)\cdot\phi\,dy
  =
  -\int_{K_1}V_n\cdot(y\cdot\nabla\phi+3\phi)\,dy,
\]
because \(\phi\) has zero trace.  Hence
\[
  |\alpha_n(s)|\,\left|\int_{K_1}(y\cdot\nabla V_n)\cdot\phi\,dy\right|
  \le C_{K_1,m}|\alpha_n(s)|\|V_n(s)\|_{L^2(K_1)}\|\phi\|_{H^m(K_1)}.
\]
Similarly,
\[
  \left|\int_{K_1}(\beta_n\cdot\nabla V_n)\cdot\phi\,dy\right|
  =
  \left|\int_{K_1}V_n\cdot(\beta_n\cdot\nabla\phi)\,dy\right|
  \le C_m|\beta_n(s)|\,|K_1|^{1/2}
  \|V_n(s)\|_{L^2(K_1)}\|\phi\|_{H^m(K_1)}.
\]
Combining the estimates and taking the supremum over
\(\|\phi\|_{H^m(K_1)}\le1\), we obtain
\[
\begin{aligned}
  \|\partial_sV_n(s)\|_{H^{-m}(K_1)}
  \le C_{K_1,m}\big(&
  \|V_n(s)\|_{L^2(K_1)}^2
  +\|\nabla V_n(s)\|_{L^2(K_1)}\\
  &+\|P_n(s)-(P_n)_{K_1}(s)\|_{L^{3/2}(K_1)}
  +( |\alpha_n(s)|+|\beta_n(s)| )\|V_n(s)\|_{L^2(K_1)}
  \big).
\end{aligned}
\]
Integrating in \(s\in(0,T)\), using Assumption~\ref{ass:selected-window-estimates}, and using
\(\alpha_n\to a_\infty\), \(\beta_n\to b_\infty\) in \(L^1\), hence uniform \(L^1\)-bounds
for \(\alpha_n,\beta_n\), proves the claimed \(L^1(0,T;H^{-m})\) bound.
\end{proof}

\begin{lemma}[Local compactness on autonomous windows]\label{lem:autonomous-window-compactness}
Assume thick autonomous windows and Assumption~\ref{ass:selected-window-estimates}.  Then after
passing to a subsequence there exist
\[
  U\in L^\infty(0,T;L^2_{\mathrm{loc}}(\mathbb R^3))
  \cap L^2(0,T;H^1_{\mathrm{loc}}(\mathbb R^3)),
  \qquad
  \Pi\in L^{3/2}_{\mathrm{loc}}((0,T)\times\mathbb R^3),
\]
such that, for every compact \(K\Subset\mathbb R^3\),
\begin{enumerate}[label=(\alph*)]
  \item
  \[
    V_n\to U
    \quad\text{strongly in }L^2(0,T;L^2(K));
  \]
  \item
  \[
    V_n\rightharpoonup U
    \quad\text{weakly in }L^2(0,T;H^1(K));
  \]
  \item
  \[
    V_n\rightharpoonup^\ast U
    \quad\text{weakly-* in }L^\infty(0,T;L^2(K));
  \]
  \item
  \[
    P_n-(P_n)_K\rightharpoonup\Pi_K
    \quad\text{weakly in }L^{3/2}((0,T)\times K),
  \]
  where the local distributions \(\nabla\Pi_K\) agree on overlaps and define
  a pressure gradient \(\nabla\Pi\) modulo functions of time.
\end{enumerate}
\end{lemma}

\begin{proof}
Fix \(K_0\Subset K_1\Subset\mathbb R^3\).  Assumption~\ref{ass:selected-window-estimates}
gives boundedness of \(V_n\) in \(L^2(0,T;H^1(K_1))\), and
Lemma~\ref{lem:window-time-regularity} gives boundedness of \(\partial_sV_n\) in
\(L^1(0,T;H^{-m}(K_1))\) for \(m\) large.  Since
\[
  H^1(K_1)\Subset L^2(K_0)\hookrightarrow H^{-m}(K_1),
\]
the Aubin--Lions--Simon compactness theorem implies, after extracting a subsequence,
\[
  V_n\to U
  \quad\text{strongly in }L^2(0,T;L^2(K_0)).
\]
The local \(L^\infty L^2\) and \(L^2H^1\) bounds also give, after extraction,
\[
  V_n\rightharpoonup^\ast U
  \quad\text{in }L^\infty(0,T;L^2(K_0)),
  \qquad
  V_n\rightharpoonup U
  \quad\text{in }L^2(0,T;H^1(K_0)).
\]
The pressure bound in Assumption~\ref{ass:selected-window-estimates} gives weak compactness of
\(P_n-(P_n)_{K_0}\) in \(L^{3/2}((0,T)\times K_0)\).  A diagonal extraction over balls
\(B_j\) gives a single subsequence for every compact \(K\).  On overlaps, two pressure limits
can differ only by a function of \(s\), because both represent the same distributional pressure
gradient in the limit equation.  Thus the gradients assemble into \(\nabla\Pi\), with \(\Pi\)
defined modulo functions of time.
\end{proof}

\begin{lemma}[Compatibility of local pressure limits]\label{lem:pressure-compatibility}
Let \(K_1,K_2\Subset\mathbb R^3\) be bounded Lipschitz domains with
nonempty connected overlap \(K_{12}:=K_1\cap K_2\).  Suppose that, after passing to a common
subsequence,
\[
  P_n-(P_n)_{K_i}\rightharpoonup \Pi_i
  \quad\text{in }L^{3/2}((0,T)\times K_i),
  \qquad i=1,2.
\]
Then there exists \(c_{12}\in L^{3/2}(0,T)\) such that
\[
  \Pi_1(y,s)-\Pi_2(y,s)=c_{12}(s)
  \quad\text{for a.e. }(y,s)\in K_{12}\times(0,T).
\]
Consequently the pressure gradient \(\nabla\Pi\) obtained from the local limits is globally
well-defined on \((0,T)\times\mathbb R^3\), with the pressure itself defined modulo functions
of time.
\end{lemma}

\begin{proof}
For each \(n\),
\[
  \bigl(P_n-(P_n)_{K_1}\bigr)-\bigl(P_n-(P_n)_{K_2}\bigr)
  =(P_n)_{K_2}(s)-(P_n)_{K_1}(s)
\]
on \(K_{12}\).  Hence its spatial gradient is zero in
\(\mathcal D'((0,T)\times K_{12})\).  Passing to the weak limit gives
\[
  \nabla_y(\Pi_1-\Pi_2)=0
  \quad\text{in }\mathcal D'((0,T)\times K_{12}).
\]
Since \(K_{12}\) is connected, the standard distributional characterization of functions with
zero spatial gradient implies that \(\Pi_1-\Pi_2=c_{12}(s)\) for some
\(c_{12}\in L^{3/2}(0,T)\).  The last assertion follows because adding a time-dependent scalar
to pressure does not change \(\nabla\Pi\).
\end{proof}

\begin{theorem}[Autonomous reduced-limit theorem]\label{thm:autonomous-limit}
Assume that a branch has thick autonomous windows and satisfies
Assumption~\ref{ass:selected-window-estimates} on the same window sequence.  Let
\((U,\Pi)\) be the limit extracted in Lemma~\ref{lem:autonomous-window-compactness}.  Then
\[
  \nabla\cdot U=0,
\]
\[
  U\not\equiv0
\]
and \((U,\Pi)\) satisfies
\begin{equation}\label{eq:autonomous-limit}
  \partial_sU+(U\cdot\nabla)U+\nabla\Pi
  =\nu\Delta U+a_\infty(U+y\cdot\nabla U)+b_\infty\cdot\nabla U,
  \qquad
  \nabla\cdot U=0,
\end{equation}
in \(\mathcal D'((0,T)\times\mathbb R^3)\).  Equivalently, for every
\(\varphi\in C_c^\infty((0,T)\times\mathbb R^3;\mathbb R^3)\),
\[
  \int_0^T\!\int
  \bigl(
  -U\cdot\partial_s\varphi
  -(U\otimes U):\nabla\varphi
  +\nu\nabla U:\nabla\varphi
  -\Pi\,\nabla\cdot\varphi
  -a_\infty(U+y\cdot\nabla U)\cdot\varphi
  -(b_\infty\cdot\nabla U)\cdot\varphi
  \bigr)\,dy\,ds=0.
\]
\end{theorem}

\begin{remark}
The theorem does not assert an autonomous Liouville or rigidity contradiction.  It proves the
compact finite-window autonomous limit from the selected-window estimates.  Any later rigidity
step must apply a separate Liouville theorem to the class containing \((U,\Pi)\).
\end{remark}

\begin{proof}
\emph{Step 1: divergence-free property of the limit.}

Each \(V_n\) is divergence free by Assumption~\ref{ass:selected-window-estimates}, so for any compactly supported scalar
\(\psi\in C_c^\infty((0,T)\times\mathbb R^3)\),
\[
  \int_{(0,T)\times\mathbb R^3}V_n\cdot\nabla\psi\,dy\,ds=0.
\]
Hence, for every such \(\psi\), since \(\nabla\psi\in C_c^\infty((0,T)\times\mathbb R^3)\) and \(V_n\to U\) in
\(L^2(0,T;L^2_{\mathrm{loc}})\),
\[
  \int_{(0,T)\times\mathbb R^3}(V_n-U)\cdot\nabla\psi\,dy\,ds\to0,
\]
we deduce
\[
  \int_{(0,T)\times\mathbb R^3}U\cdot\nabla\psi\,dy\,ds=0.
\]
Hence \(\nabla\cdot U=0\) in \(\mathcal D'((0,T)\times\mathbb R^3)\).


\emph{Step 2: nontriviality of the limit profile.}
Define for \(s\in(0,T)\)
\[
  f_n(s):=\int_{\mathbb R^3}|V_n(y,s)|^2\chi(y)
  \,dy,
  \qquad
  f(s):=\int_{\mathbb R^3}|U(y,s)|^2\chi(y)
  \,dy.
\]
Here \(K_\chi:=\operatorname{supp}\chi\), so
\(\chi^{1/2}V_n,\chi^{1/2}U\in L^2(0,T;L^2(K_\chi))\).
Then
\[
  f_n-f=\int_{\mathbb R^3}(V_n-U)\cdot(V_n+U)\chi
  \,dy.
\]
Applying Cauchy--Schwarz in space,
\[
  |f_n(s)-f(s)|\le
  \|\chi^{1/2}(V_n(s)-U(s))\|_{L^2(\mathbb R^3)}
  \|\chi^{1/2}(V_n(s)+U(s))\|_{L^2(\mathbb R^3)}.
\]
Integrating in \(s\) and applying Cauchy--Schwarz in time gives
\[
  \|f_n-f\|_{L^1(0,T)}
  \le
  \|\chi^{1/2}(V_n-U)\|_{L^2(0,T;L^2)}
  \|\chi^{1/2}(V_n+U)\|_{L^2(0,T;L^2)}.
\]
Since \(K_\chi\) is compact, Assumption~\ref{ass:selected-window-estimates} and
the membership \(U\in L^2(0,T;L^2(K_\chi))\) give
\[
  \sup_n\|\chi^{1/2}(V_n+U)\|_{L^2(0,T;L^2)}
  \le C_\chi\left(\sup_n\|V_n\|_{L^2(0,T;L^2(K_\chi))}
  +\|U\|_{L^2(0,T;L^2(K_\chi))}\right)<\infty.
\]
Moreover
\[
  \|\chi^{1/2}(V_n-U)\|_{L^2(0,T;L^2)}
  \le \|\chi\|_{L^\infty}^{1/2}
  \|V_n-U\|_{L^2(0,T;L^2(K_\chi))}\to0.
\]
Therefore
\[
  \|f_n-f\|_{L^1(0,T)}\to0.
\]
Hence \(f_n\to f\) in \(L^1(0,T)\).

Assumption~\ref{ass:selected-window-estimates} gives for large \(n\),
\(f_n(s)\ge\eta\) for a.e.\(s\). Therefore
\[
  \int_0^T f(s)
  \,ds=\lim_{n\to\infty}\int_0^Tf_n(s)
  \,ds\ge\eta T>0.
\]
Hence \(f\not\equiv0\), so \(U\not\equiv0\).

\emph{Step 3: shifted weak formulation for \(V_n\).}
Take any test field
\(\varphi\in C_c^\infty((0,T)\times\mathbb R^3;\mathbb R^3)\), and set
\[
K_\varphi:=\overline{\{x\in\mathbb R^3:\exists s\in(0,T),\ \varphi(x,s)\neq0\}}
\subset\mathbb R^3,\qquad K_\varphi\ \text{bounded}.
\]
Since \((V,P)\) satisfies~\eqref{eq:renorm-ns} in distributions, the change of variables
\(\tau=\tau_n+s\) gives the shifted weak formulation
\begin{equation}\label{eq:weak-n}
\begin{aligned}
0={}&\int_0^T\!\int_{\mathbb R^3}
\Bigl(
-V_n\cdot\partial_s\varphi
-(V_n\otimes V_n):\nabla\varphi
+\nu\nabla V_n:\nabla\varphi
-P_n\nabla\cdot\varphi
\Bigr)\,dy\,ds\\
&+\int_0^T\!\int_{\mathbb R^3}
\Bigl(
-a(\tau_n+s)(V_n+y\cdot\nabla V_n)\cdot\varphi
-(b(\tau_n+s)\cdot\nabla V_n)\cdot\varphi
\Bigr)\,dy\,ds.
\end{aligned}
\end{equation}

\emph{Step 4: term-by-term passage to the limit.}
Write \(\alpha_n(s):=a(\tau_n+s)\), \(\beta_n(s):=b(\tau_n+s)\).

\paragraph{(i) Time derivative.}

By strong convergence in \(L^2(0,T;L^2(K_\varphi))\),
\[
  \left|\int_0^T\!\int (V_n-U)\cdot\partial_s\varphi
  \,dy\,ds\right|
  \le
  \|V_n-U\|_{L^2(0,T;L^2(K_\varphi))}
  \|\partial_s\varphi\|_{L^2(0,T;L^2(K_\varphi))}
  \to0.
\]
Hence
\[\int V_n\cdot\partial_s\varphi\to\int U\cdot\partial_s\varphi.
\]

\paragraph{(ii) Convection term.}
Decompose
\[
  V_n\otimes V_n-U\otimes U=(V_n-U)\otimes V_n+U\otimes(V_n-U).
\]
Then
\[
\begin{aligned}
&\left|\int_0^T\!\int((V_n\otimes V_n-U\otimes U):\nabla\varphi)
  \,dy\,ds\right|\\
&\quad\le
  \|\nabla\varphi\|_{L^\infty_{x,t}}
  \|V_n-U\|_{L^2(0,T;L^2(K_\varphi))}
  \bigl(\|V_n\|_{L^2(0,T;L^2(K_\varphi))}
  +\|U\|_{L^2(0,T;L^2(K_\varphi))}\bigr)
  \to0,
\end{aligned}
\]
using boundedness of \(V_n\) in \(L^2(0,T;L^2(K_\varphi))\). Therefore
\(
  \int(V_n\otimes V_n):\nabla\varphi
  \to\int(U\otimes U):\nabla\varphi.
\)

\paragraph{(iii) Viscous term.}
From weak convergence of \(V_n\) in \(L^2(0,T;H^1_{\mathrm{loc}})\),
\[
  \int_0^T\!\int\nabla V_n:\nabla\varphi
  \,dy\,ds
  \to
  \int_0^T\!\int\nabla U:\nabla\varphi
  \,dy\,ds.
\]

\paragraph{(iv) Pressure term.}
Let \(K\) be a bounded Lipschitz domain with \(K_\varphi\Subset K\).  Since
\(\varphi\) is compactly supported in \(K\),
\[
  \int_K (P_n)_K(s)\,\nabla\cdot\varphi(y,s)\,dy
  =
  (P_n)_K(s)\int_K\nabla\cdot\varphi(y,s)\,dy=0.
\]
Therefore
\[
  \int_0^T\!\int_{\mathbb R^3}P_n\nabla\cdot\varphi\,dy\,ds
  =
  \int_0^T\!\int_K(P_n-(P_n)_K)\nabla\cdot\varphi\,dy\,ds.
\]
By Lemma~\ref{lem:autonomous-window-compactness},
\[
  P_n-(P_n)_K\rightharpoonup\Pi_K
  \quad\text{in }L^{3/2}((0,T)\times K),
\]
and \(\nabla\cdot\varphi\in L^3((0,T)\times K)\).  The compatibility
Lemma~\ref{lem:pressure-compatibility} permits us to write the resulting pressure distribution
as \(\Pi\), modulo functions of time.  Hence
\[
  \int_0^T\!\int_{\mathbb R^3}P_n\nabla\cdot\varphi\,dy\,ds
  \to
  \int_0^T\!\int_K\Pi_K\nabla\cdot\varphi\,dy\,ds.
\]
The value of the last integral is independent of the chosen \(K\), because replacing
\(\Pi_K\) by \(\Pi_K+c(s)\) changes the integral by
\(\int_0^Tc(s)\int_K\nabla\cdot\varphi\,dy\,ds=0\).  This is the pressure term
\(\int\Pi\nabla\cdot\varphi\).

\paragraph{(v) Drift term with \(a\): part involving \(V_n\).}
Define
\[
  Y_n^0(s):=\int_{\mathbb R^3}V_n(y,s)\cdot\varphi(y,s)
  \,dy,
  \qquad
  Y^0(s):=\int_{\mathbb R^3}U(y,s)\cdot\varphi(y,s)
  \,dy.
\]
For each \(s\), Cauchy--Schwarz gives
\[
  |Y_n^0(s)|
  \le
  \|V_n(s)\|_{L^2(K_\varphi)}
  \|\varphi(s)\|_{L^2(K_\varphi)}.
\]
By Assumption~\ref{ass:selected-window-estimates}, there is \(C_0<\infty\)
such that \(\sup_n\|Y_n^0\|_{L^\infty(0,T)}\le C_0\). Moreover,
\[
  Y_n^0\to Y^0\quad\text{in }L^2(0,T),
\]
because \(V_n\to U\) in \(L^2(0,T;L^2(K_\varphi))\).
Now
\[
  \int_0^T\alpha_n(s)Y_n^0(s)
  \,ds
  =\int_0^T(\alpha_n(s)-a_\infty)Y_n^0(s)
  \,ds+a_\infty\int_0^TY_n^0(s)
  \,ds.
\]
For the first term,
\[
  \left|\int_0^T(\alpha_n-a_\infty)Y_n^0\,ds\right|\le
\sup_{s\in(0,T)}|Y_n^0(s)|\,\|\alpha_n-a_\infty\|_{L^1(0,T)}\to0.
\]
For the second term, strong convergence in \(L^2(0,T)\) implies
\(\int_0^TY_n^0\to\int_0^TY^0\), so
\[
  \int_0^T\alpha_n(s)\int_{\mathbb R^3}V_n\cdot\varphi
  \,dy\,ds
  \to a_\infty\int_0^T\int_{\mathbb R^3}U\cdot\varphi
  \,dy\,ds.
\]
Thus the signed contribution in \eqref{eq:weak-n} satisfies
\[
  -\int_0^T\alpha_n(s)\int_{\mathbb R^3}V_n\cdot\varphi
  \,dy\,ds
  \to -a_\infty\int_0^T\int_{\mathbb R^3}U\cdot\varphi
  \,dy\,ds.
\]

\paragraph{(vi) Drift term with \(a\): part involving \(y\cdot\nabla V_n\).}
Since \(\varphi\in C_c^\infty\), integrating by parts in \(y\) gives
\[
  \int_{\mathbb R^3}(y\cdot\nabla V_n)\cdot\varphi
  =-\int_{\mathbb R^3}V_n\cdot(y\cdot\nabla\varphi+3\varphi)
  \,dy.
\]
Define
\[
  Y_n^1(s):=\int_{\mathbb R^3}V_n\cdot(y\cdot\nabla\varphi+3\varphi)
  \,dy,
  \qquad
  Y^1(s):=\int_{\mathbb R^3}U\cdot(y\cdot\nabla\varphi+3\varphi)
  \,dy.
\]
For each \(s\), Cauchy--Schwarz gives
\[
  |Y_n^1(s)|
  \le
  \|V_n(s)\|_{L^2(K_\varphi)}
  \|y\cdot\nabla\varphi(s)+3\varphi(s)\|_{L^2(K_\varphi)}.
\]
The \(L^\infty(0,T;L^2(K_\varphi))\) bound in
Assumption~\ref{ass:selected-window-estimates} gives
\[
  \sup_n\|Y_n^1\|_{L^\infty(0,T)}\le C_1<\infty.
\]
Also,
\[
  \|Y_n^1-Y^1\|_{L^2(0,T)}
  \le
  \|V_n-U\|_{L^2(0,T;L^2(K_\varphi))}
  \|y\cdot\nabla\varphi+3\varphi\|_{L^\infty(0,T;L^2(K_\varphi))}
  \to0.
\]
Therefore
\[
  -\int_0^T\alpha_n(s)
  \int_{\mathbb R^3}(y\cdot\nabla V_n)\cdot\varphi
  \,dy\,ds
  \to
  a_\infty\int_0^T\!\int_{\mathbb R^3}U\cdot(y\cdot\nabla\varphi+3\varphi)
  \,dy\,ds.
\]
Integrating by parts once more in the limit gives
\[
  a_\infty\int_0^T\!\int_{\mathbb R^3}U\cdot(y\cdot\nabla\varphi+3\varphi)
  \,dy\,ds
  =
  -a_\infty\int_0^T\!\int_{\mathbb R^3}(y\cdot\nabla U)\cdot\varphi
  \,dy\,ds.
\]

\paragraph{(vii) Term with \(b\).}
Since \(\beta_n(s)\) is independent of \(y\),
\[
  \int_{\mathbb R^3}(\beta_n\cdot\nabla V_n)\cdot\varphi
  =-\int_{\mathbb R^3}V_n\cdot(\beta_n\cdot\nabla\varphi)
  \,dy,
\]
so
\[
  -\int_0^T\!\int(\beta_n\cdot\nabla V_n)\cdot\varphi
  \,dy\,ds
  =
  \int_0^T\!\int V_n\cdot(\beta_n\cdot\nabla\varphi)
  \,dy\,ds.
\]
Let
\[
\begin{aligned}
Z_n(s)&:=\left(\int_{\mathbb R^3}V_n\cdot\partial_1\varphi\,dy,\,
\int_{\mathbb R^3}V_n\cdot\partial_2\varphi\,dy,\,
\int_{\mathbb R^3}V_n\cdot\partial_3\varphi\,dy\right),\\
Z(s)&:=\left(\int_{\mathbb R^3}U\cdot\partial_1\varphi\,dy,\,
\int_{\mathbb R^3}U\cdot\partial_2\varphi\,dy,\,
\int_{\mathbb R^3}U\cdot\partial_3\varphi\,dy\right).
\end{aligned}
\]
Then
\[
  \int_{\mathbb R^3}V_n\cdot(\beta_n\cdot\nabla\varphi)\,dy
  =\beta_n(s)\cdot Z_n(s),\qquad
  \int_{\mathbb R^3}U\cdot(b_\infty\cdot\nabla\varphi)\,dy
  =b_\infty\cdot Z(s).
\]
Consequently
\[
  \left|\int_0^T\!(\beta_n-b_\infty)\cdot Z_n\,ds\right|
  \le
  \|\beta_n-b_\infty\|_{L^1(0,T)}\sup_{s\in(0,T)}|Z_n(s)|\to0,
\]
because
\[
  |Z_n(s)|\le
  \|V_n(s)\|_{L^2(K_\varphi)}
  \left(\sum_{j=1}^3\|\partial_j\varphi(s)\|_{L^2(K_\varphi)}^2\right)^{1/2}
\]
and Assumption~\ref{ass:selected-window-estimates} bounds the right-hand side
uniformly in \(n\) and \(s\). Furthermore,
\[
  \|Z_n-Z\|_{L^2(0,T;\mathbb R^3)}
  \le
  \|V_n-U\|_{L^2(0,T;L^2(K_\varphi))}
  \left(\sum_{j=1}^3\|\partial_j\varphi\|_{L^\infty(0,T;L^2(K_\varphi))}^2\right)^{1/2}
  \to0.
\]
Therefore
\[
  \int_0^T b_\infty\cdot Z_n(s)\,ds
  \to
  \int_0^T b_\infty\cdot Z(s)\,ds.
\]
Hence
\[
  -\int_0^T\!\int (\beta_n\cdot\nabla V_n)\cdot\varphi
  \,dy\,ds
  \to
  \int_0^T\!\int U\cdot(b_\infty\cdot\nabla\varphi)
  \,dy\,ds.
\]
Since \(b_\infty\) is constant in \(y\), integration by parts yields
\[
  \int_0^T\!\int U\cdot(b_\infty\cdot\nabla\varphi)
  \,dy\,ds
  =
  -\int_0^T\!\int (b_\infty\cdot\nabla U)\cdot\varphi
  \,dy\,ds.
\]

\emph{Step 5: assemble the limit.}
Passing to the limit in \eqref{eq:weak-n} in each term (Steps 4(i)--4(vii)) gives
\[
  \int_0^T\!\int
  \bigl(
  -U\cdot\partial_s\varphi
  -(U\otimes U):\nabla\varphi
  +\nu\nabla U:\nabla\varphi
  -\Pi\nabla\cdot\varphi
  -a_\infty(U+y\cdot\nabla U)\cdot\varphi
  -(b_\infty\cdot\nabla U)\cdot\varphi
  \bigr)
  \,dy\,ds=0.
\]
Since this holds for every compactly supported vector field \(\varphi\), equation~\eqref{eq:autonomous-limit} holds in distributions on
\((0,T)\times\mathbb R^3\).
\end{proof}

\section{State-space stratification of the scale-collapse branch}
\label{sec:branch-closure}

The cost argument above closes the finite-cost scale-collapse branch.  The remaining
scale-collapse obstruction is treated by a state-space stratification, following the local
estimate and autonomous-reduction route of the preceding papers.  The purpose of this section is
not to assert a blanket Liouville theorem for every autonomous negative-drift equation.  Instead,
we separate the closed branches from the precise residual defects.

\begin{theorem}[Selected core retention]\label{thm:core-retention}
Let a represented Type II branch be obtained from the selected-core construction of the preceding
papers.  Then there exist \(R>0\), a cutoff \(\chi\in C_c^\infty(B_R)\), \(\chi\ge0\), a
constant \(\eta>0\), and a terminal time \(\tau_1\) such that, along every selected window used
below,
\[
  \int_{\mathbb R^3}|V(y,\tau)|^2\chi(y)\,dy\ge\eta
  \quad\text{for a.e. }\tau\ge\tau_1.
\]
In particular, every translated sequence \(V_n(y,s)=V(y,\tau_n+s)\) satisfies the localized
nontriviality condition in Assumption~\ref{ass:selected-window-estimates}, after discarding
finitely many \(n\).
\end{theorem}

\begin{proof}
The selected-core construction chooses a compact renormalized core carrying a fixed positive
amount of localized kinetic mass.  Since the cutoff \(\chi\) is supported strictly inside that
core and equals a positive spatial weight there, the selected-core lower bound gives the stated
estimate.  Translating \(\tau=\tau_n+s\) gives the final assertion on every finite selected
window.
\end{proof}

\begin{theorem}[Autonomous modulation selection]\label{thm:autonomous-modulation-selection}
Let a represented scale-collapse branch satisfy the modulation bounds obtained in the repaired
gauge construction.  Suppose that a sequence of windows \([\tau_n,\tau_n+T]\) is selected in the
autonomous stratum.  Then, after passing to a subsequence, there exist constants
\(a_\infty\in\mathbb R\) and \(b_\infty\in\mathbb R^3\) such that
\[
  a(\tau_n+\cdot)\to a_\infty
  \quad\text{in }L^1(0,T),
  \qquad
  b(\tau_n+\cdot)\to b_\infty
  \quad\text{in }L^1(0,T;\mathbb R^3).
\]
\end{theorem}

\begin{proof}
This is the autonomous-window selection conclusion of the repaired-gauge modulation analysis:
the selected stratum is defined by convergence of the modulation functions to constant
autonomous parameters on the fixed window length \(T\).  Passing to the subsequence supplied by
that selection gives exactly the displayed \(L^1\)-convergences.
\end{proof}

\begin{theorem}[Scale-collapse state-space stratification]\label{thm:scale-collapse-stratification}
Let a represented Type II branch have genuine scale collapse and nonvanishing selected core.
Then exactly one of the following alternatives occurs after passing to selected terminal
windows:
\begin{enumerate}[label=(\roman*)]
  \item the branch has finite scale-collapsing cost, which is excluded by
  Theorem~\ref{thm:cost-exclusion};
  \item the branch has finite absolute cost, which is excluded by
  Corollary~\ref{cor:absolute-cost-exclusion};
  \item the branch admits thick autonomous windows in the sense of
  Definition~\ref{def:thick-autonomous-windows}, satisfies the selected-window estimates of
  Assumption~\ref{ass:selected-window-estimates}, and has autonomous modulation convergence;
  \item the branch has a thin-drift defect: no selected fixed-length window carries a uniform
  positive lower bound for the averaged weighted negative drift;
  \item the branch has an autonomous-modulation defect: thick windows exist, but no subsequence
  has a constant-coefficient \(L^1\) modulation limit in the sense needed for
  Theorem~\ref{thm:autonomous-limit};
  \item the branch has a compactness defect: the local estimates needed for
  Assumption~\ref{ass:selected-window-estimates} fail on the selected windows.
\end{enumerate}
\end{theorem}

\begin{proof}
The scale-collapse selection theorem decomposes terminal scale-collapse windows according to the
averaged negative-drift mass
\[
  \frac1T\int_{\tau_n}^{\tau_n+T}a_-(\tau)M_R(\tau)\,d\tau.
\]
If the accumulated cost is finite, alternatives (i) or (ii) apply according to the cost
functional under consideration.  If the averaged negative-drift mass stays bounded below on a
subsequence, one obtains thick windows; otherwise the branch is in the thin-drift alternative.
On thick windows, either the local estimates and core retention hold, in which case
Proposition~\ref{prop:selected-window-derived} and Theorem~\ref{thm:core-retention} give
Assumption~\ref{ass:selected-window-estimates}, or the branch is in the compactness defect.  If
the estimates hold but no subsequence has the required constant-coefficient modulation limit,
the branch is in the autonomous-modulation defect.  The remaining case is precisely the thick
autonomous alternative.
\end{proof}

\begin{definition}[Stationary omega-limit class]\label{def:stationary-omega-class}
Let \((U,\Pi)\) be a nonzero autonomous limit produced by
Theorem~\ref{thm:autonomous-limit}.  We say that it has a stationary omega-limit in the
admissible class if there exists a nonzero pair
\[
  (W,Q),\qquad W\in L^3(\mathbb R^3),
\]
with the local energy and pressure reconstruction properties inherited from the selected
windows, such that
\[
  (W\cdot\nabla)W+\nabla Q
  =
  \nu\Delta W+a_\infty(W+y\cdot\nabla W)+b_\infty\cdot\nabla W,
  \qquad
  \nabla\cdot W=0,
\]
in \(\mathcal D'(\mathbb R^3)\).
\end{definition}

\begin{definition}[Rigidity-covered stationary class]\label{def:rigidity-covered-class}
For a parameter pair \((a_\infty,b_\infty)\), the stationary class is rigidity-covered if every
admissible stationary solution in Definition~\ref{def:stationary-omega-class} is either zero or
is not a terminal Type II branch.  In the classical backward self-similar normalization this is
the role of the Nečas--Růžička--Šverák rigidity theorem, after the harmless translation drift is
removed.  Other parameter regimes require their own stated Liouville or classification theorem.
\end{definition}

\begin{theorem}[Stationary-rigidity blocker]\label{thm:stationary-rigidity-blocker}
Assume a thick autonomous branch produces a nonzero autonomous limit by
Theorem~\ref{thm:autonomous-limit}.  If that autonomous limit has a stationary omega-limit in
the admissible class and the corresponding stationary class is rigidity-covered, then the branch
is not an unresolved Type II scale-collapse branch.
\end{theorem}

\begin{proof}
The stationary omega-limit gives a nonzero stationary profile \((W,Q)\) in the admissible class.
The rigidity-covered hypothesis says that no such nonzero profile can represent a terminal Type
II branch.  This contradicts unresolved membership in the scale-collapse Type II state.
\end{proof}

\begin{theorem}[Scale-collapse alternatives]\label{thm:scale-collapse-alternatives}
Let a represented Type II branch have genuine scale collapse and a nonvanishing selected core.
Then one of the following mutually exclusive outcomes occurs:
\begin{enumerate}[label=(\roman*)]
  \item the finite scale-collapsing cost exclusion applies;
  \item the finite absolute cost exclusion applies;
  \item the branch has a thin-drift defect;
  \item the branch has an autonomous-modulation defect;
  \item the branch has a compactness defect;
  \item the branch has a nonzero autonomous reduced limit, and either a stationary-rigidity
  blocker applies or the remaining obstruction is the corresponding stationary/generalized
  self-similar Liouville problem.
\end{enumerate}
\end{theorem}

\begin{proof}
Apply Theorem~\ref{thm:scale-collapse-stratification}.  The finite-cost alternatives are exactly
items (i) and (ii).  The thin, modulation, and compactness defects are items (iii)--(v).  In the
remaining thick autonomous case, Lemma~\ref{lem:negative-parameter} gives \(a_\infty<0\), and
Theorem~\ref{thm:autonomous-limit} produces a nonzero autonomous reduced limit.  If
Definitions~\ref{def:stationary-omega-class} and~\ref{def:rigidity-covered-class} apply, then
Theorem~\ref{thm:stationary-rigidity-blocker} blocks the branch.  If not, the unresolved
remainder is precisely the stationary or generalized self-similar Liouville problem for the
parameter pair \((a_\infty,b_\infty)\).
\end{proof}

\section{Conclusion}

The exclusions in Sections~\ref{sec:cost}--\ref{sec:finite-cost} are proved unconditionally within
the renormalized setting:
no branch with finite scale-collapsing (hence finite absolute) cost can have
\(\lambda(\tau)\to0\) together with a nonvanishing localized core.

Lemma~\ref{lem:negative-parameter} shows that any thick autonomous window admits a negative
autonomous drift limit \(a_\infty\).  Lemmas~\ref{lem:window-time-regularity}
and~\ref{lem:autonomous-window-compactness}, together with Theorem~\ref{thm:autonomous-limit},
assemble the selected-window estimates from the preceding papers into a nontrivial autonomous
finite-window limit \((U,\Pi)\) satisfying the full distributional autonomous equation.

Section~\ref{sec:branch-closure} records the same state-space stratification as the markdown
route.  The autonomous reduction is a closed contradiction only in parameter regimes where a
stationary omega-limit is available and the resulting stationary class is covered by a Liouville
or rigidity theorem, such as the classical Nečas--Růžička--Šverák self-similar class.  Outside
those regimes the paper records the precise thin-drift, modulation, compactness, or
Liouville-defect alternative rather than asserting an unconditional closure.

\end{document}
